\documentclass[11pt,a4paper,reqno]{amsart}

\usepackage[utf8]{inputenc}
\usepackage[T1]{fontenc}
\usepackage{amsmath,amssymb,amsfonts,amsthm,mathtools}
\usepackage{mathrsfs}
\usepackage{microtype}
\usepackage{booktabs}
\usepackage{tabularx}
\usepackage[dvipsnames]{xcolor}
\usepackage{hyperref}
\usepackage{geometry}
\usepackage{enumitem}

\geometry{
  top=3cm,
  bottom=3cm,
  left=2.8cm,
  right=2.8cm,
  headheight=14pt
}

\hypersetup{
  colorlinks=true,
  linkcolor=NavyBlue,
  citecolor=Mahogany,
  urlcolor=TealBlue,
  pdfauthor={Reinaldo M. Silva-Filho},
  pdftitle={Beyond the Spectrum II: Metric Measure Geometry, Persistent Homology, and Non-Commutative Invariants of Functional Tensor Manifolds}
}

% --- Theorem Environments ---
\newtheorem{theorem}{Theorem}[section]
\newtheorem{lemma}[theorem]{Lemma}
\newtheorem{proposition}[theorem]{Proposition}
\newtheorem{corollary}[theorem]{Corollary}

\theoremstyle{definition}
\newtheorem{definition}[theorem]{Definition}
\newtheorem{example}[theorem]{Example}
\newtheorem{remark}[theorem]{Remark}
\newtheorem{construction}[theorem]{Construction}

\numberwithin{equation}{section}

% --- Custom Notations ---
\newcommand{\R}{\mathbb{R}}
\newcommand{\C}{\mathbb{C}}
\newcommand{\N}{\mathbb{N}}
\newcommand{\Z}{\mathbb{Z}}
\newcommand{\Sph}{\mathbb{S}}
\newcommand{\Tor}{\mathbb{T}}
\newcommand{\Id}{\mathrm{Id}}
\newcommand{\Tr}{\operatorname{Tr}}
\newcommand{\rank}{\operatorname{rank}}
\newcommand{\supp}{\operatorname{supp}}
\newcommand{\Lip}{\operatorname{Lip}}
\newcommand{\BV}{\operatorname{BV}}
\newcommand{\TV}{\operatorname{TV}}
\newcommand{\Per}{\operatorname{Per}}
\newcommand{\cutnorm}[1]{\|#1\|_{\square}}
\newcommand{\norm}[1]{\left\|#1\right\|}
\newcommand{\inner}[2]{\left\langle #1, #2 \right\rangle}
\newcommand{\dif}{\,\mathrm{d}}
\newcommand{\abs}[1]{\left|#1\right|}
\newcommand{\WF}{\operatorname{WF}}
\newcommand{\QFI}{\mathrm{QFI}}
\newcommand{\Ric}{\operatorname{Ric}}
\newcommand{\Hess}{\operatorname{Hess}}
\newcommand{\diam}{\operatorname{diam}}
\newcommand{\dist}{\operatorname{dist}}

\begin{document}

\title[Beyond the Spectrum II]{Beyond the Spectrum II: Metric Measure Geometry, Persistent Homology, and Non-Commutative Invariants of Functional Tensor Manifolds}

\author{Reinaldo M. Silva-Filho}
\address{Universidade Federal de Lavras (UFLA), Departamento de Estat\'istica e Experimenta\c{c}\~ao Agropecu\'aria (PPGEEAA/DES)}
\email{reinaldo.filho1@estudante.ufla.br}
\date{\today}

\begin{abstract}
Finite-dimensional multilinear algebra classically relies on permutation-blind algebraic invariants: eigenvalues, singular values, operator ranks, and polynomial traces. While the initial formulation of \emph{Functional Realization Operators} $\Phi \colon \mathcal{T}^k(V) \to \mathcal{F}(\Omega)$ demonstrated that embedding matrices and higher-order tensors into function spaces unlocks first-order Sobolev $H^1$ energies, total variation bounded variation ($\BV$) norms, and graphon cut-norm topologies, profound geometric and topological dimensions have remained unchartered. In this paper, we establish a comprehensive higher-order framework by constructing six new foundational families of emergent invariants and geometric measures on functional tensor varieties:
(1) \textbf{Optimal Transport and Metric Curvature:} a dimension-free quadratic Wasserstein distance $\mathcal{W}_2$ measuring continuous deformation costs between matrices of unequal dimensions, paired with a Bakry--\'Emery Ricci curvature lower bound $\kappa_{\mathrm{BE}}$ governing displacement convexity and logarithmic Sobolev inequalities;
(2) \textbf{Topological Data Analysis and Persistent Homology:} continuous super-level set filtrations $X_t(A)$ yielding persistence diagrams $\mathrm{Dgm}_k(\Phi(A))$, proving strict bottleneck stability $d_B \le \|\Phi(A) - \Phi(B)\|_\infty$ and establishing persistent entropy as an invariant that strictly separates isospectral matrices;
(3) \textbf{Geometric Measure Theory and Bending Energies:} level-set varifolds defining the integrated Willmore bending energy $\mathcal{W}(\Phi(A))$ to quantify surface wrinkling and mean-curvature tortuosity under conformal transformations, alongside matrix isoperimetric profiles;
(4) \textbf{Micro-Local Phase Space and Fractional Regularity:} continuous Gabor--Radon wavepacket transforms defining the matrix wavefront set $\mathrm{WF}(\Phi(A)) \subset T^*\Omega$ to detect the spatial location and directional propagation of block discontinuities, paired with critical fractional Besov exponents $s^*(A)$;
(5) \textbf{Non-Commutative Geometry:} canonical Dirac spectral triples $(\mathcal{A}_\Phi, L^2(\Omega, \mathbb{S}), \mathcal{D})$ where the Connes--Dixmier trace $\Tr_\omega(\Phi(A)|\mathcal{D}|^{-d})$ extracts conformal volume elements, alongside quantized topological cyclic cocycles $\tau_{2k}$;
(6) \textbf{Quantum Information Geometry of Tensor Networks:} fiber-bundle continuous matrix product states defining the integrated Quantum Fisher Information (QFI) volume $\mathrm{Vol}_{\mathrm{QFI}}(\mathcal{T})$, continuous entanglement contours $\mathcal{S}_{\mathrm{cont}}(x)$, and spatial information centroids.
We prove rigorous existence, stability, and invariance theorems for each construction, providing an exhaustive classification that bridges continuous multilinear algebra, global differential geometry, and theoretical physics.
\end{abstract}

\keywords{Functional linear algebra, higher-order tensors, optimal transport, Wasserstein distance, Bakry--\'Emery curvature, persistent homology, topological data analysis, Willmore energy, geometric measure theory, wavefront set, micro-local analysis, non-commutative geometry, Dixmier trace, Quantum Fisher Information, tensor networks.}
\subjclass[2020]{15A69, 49Q22, 55N31, 53C21, 58J40, 58B34, 81P45, 46E35}

\maketitle

\tableofcontents

\newpage

% =========================================================================
\section{Introduction and the Functional Realization Paradigm}
\label{sec:intro}
% =========================================================================

Classical linear and multilinear algebra views an array of numerical coefficients---be it a square matrix $A \in \R^{n \times n}$ or a higher-order tensor $\mathcal{T} \in \bigotimes_{j=1}^k \R^{d_j}$---as an element of a finite-dimensional discrete vector space. In this classical framework, all intrinsic structural characterizations are governed by purely algebraic invariants: the matrix trace $\Tr(A)$, the determinant $\det(A)$, characteristic polynomials, the eigenvalue spectrum $\sigma(A) = \{\lambda_i(A)\}_{i=1}^n$, singular value decompositions $\Sigma(A)$, and multilinear ranks.

However, as established in the foundational treatise \cite{silvafilho2026beyond}, this algebraic perspective suffers from severe structural blind spots when arrays represent physical, geometric, or network phenomena:
\begin{enumerate}[label=(\roman*)]
    \item \textbf{Spectral Blindness to Spatial Locality:} Algebraic invariants are invariant under coordinate permutations. If $P \in \R^{n \times n}$ is a permutation matrix, then
    \begin{equation}
        \sigma(PAP^T) = \sigma(A), \quad \Tr(PAP^T) = \Tr(A), \quad \|PAP^T\|_F = \|A\|_F.
    \end{equation}
    Yet, transposing adjacent rows versus distant rows in an image or adjacency matrix creates radically different spatial gradients, topological tears, and boundary discontinuities.
    \item \textbf{Dimensional Incommensurability:} Two matrices of different dimensions, say $A \in \R^{10 \times 10}$ and $B \in \R^{1000 \times 1000}$, inhabit non-isomorphic linear spaces. Classical matrix algebra possesses no natural, canonical metric to measure how closely $A$ and $B$ represent the same continuous physical continuum.
\end{enumerate}

To overcome these limitations, the framework of \emph{Functional Realization Operators} was introduced:
\begin{equation}
    \Phi \colon \mathcal{T}^k(V) \longrightarrow \mathcal{F}(\Omega),
\end{equation}
mapping discrete multilinear arrays into function spaces (such as $L^p(\Omega)$, Sobolev spaces $H^s(\Omega)$, or spaces of bounded variation $\BV(\Omega)$) over continuous spatial domains $\Omega \subset \R^d$ or Riemannian manifolds $(\Omega, g)$. 

While the first volume \cite{silvafilho2026beyond} successfully explored first-order Sobolev Dirichlet energies $\int_\Omega \|\nabla \Phi(A)\|^2 \dif x$, $\BV$ level-set perimeters via the Federer--Fleming Coarea formula, spherical Morse topology for symmetric tensors, and graphon cut-norm limits, the true depth of functional realizations remained unchartered. In particular, first-order differential invariants cannot capture:
\begin{itemize}
    \item The optimal dynamical transport cost between differing mass distributions of unequal matrices;
    \item The multiscale topological persistence of voids and cavities in the entry landscape;
    \item The second-order bending, mean-curvature wrinkling, and Willmore energy of matrix level sets;
    \item The exact micro-local directional orientation of shock interfaces and block boundaries in cotangent space;
    \item Non-commutative operator traces and cyclic topological cocycles;
    \item The fiber-bundle information geometry of continuous tensor network manifolds.
\end{itemize}

In this paper, we construct and prove the mathematical foundations of these six new advanced families of invariants and geometric measures.

---

\subsection{Outline of the Treatise}
The paper is organized as follows:
\begin{itemize}
    \item \textbf{Section \ref{sec:optimal_transport}:} We construct the continuous Wasserstein metric measure space for functional tensors, establishing the dynamical $\mathcal{W}_2$ geodesic distance between matrices of arbitrary dimensions, followed by the Bakry--\'Emery $\Gamma_2$ Ricci curvature bound and logarithmic Sobolev inequalities.
    \item \textbf{Section \ref{sec:persistent_homology}:} We define the continuous super-level set filtration, prove the Bottleneck Stability Theorem for functional realizations, and construct the Persistent Entropy invariant that strictly separates isospectral matrices.
    \item \textbf{Section \ref{sec:willmore_gmt}:} We develop the Geometric Measure Theory (GMT) of level-set varifolds, formulating the Integrated Willmore Bending Energy $\mathcal{W}(\Phi(A))$ and the matrix isoperimetric profile.
    \item \textbf{Section \ref{sec:microlocal}:} We introduce continuous Gabor wavepackets on phase space, formalize the Matrix Wavefront Set $\WF(\Phi(A)) \subset T^*\Omega \setminus \{0\}$, and derive critical fractional Besov/Gagliardo regularity exponents $s^*(A)$.
    \item \textbf{Section \ref{sec:noncommutative}:} We construct the canonical Dirac spectral triple, establish the Connes--Dixmier non-commutative volume trace $\Tr_\omega$, and formalize quantized cyclic cocycles $\tau_{2k}$.
    \item \textbf{Section \ref{sec:quantum_info}:} We formulate fiber-bundle tensor networks, the integrated Quantum Fisher Information (QFI) volume, the continuous entanglement contour, and the spatial information centroid.
    \item \textbf{Section \ref{sec:synthesis}:} We present the unified master comparison matrix, computational discrete-to-continuous realization algorithms, and applications across deep learning, disordered materials, and quantum information.
\end{itemize}

% =========================================================================
\section{Optimal Transport, Wasserstein Geodesics, and Metric Curvature}
\label{sec:optimal_transport}
% =========================================================================

\subsection{The Density Realization Operator}
Let $V$ be an $n$-dimensional Euclidean space. Let $\mathcal{S}_+^n \subset \R^{n \times n}$ denote the cone of symmetric positive semi-definite (PSD) matrices, normalized such that $\Tr(A) = 1$. Let $\Omega = [0, 1]^d$ be the unit hypercube equipped with the standard Lebesgue measure $\dif x$.

\begin{definition}[Probability Density Realization]
Let $\Phi \colon \mathcal{S}_+^n \to L^1(\Omega)$ be a positive functional realization operator satisfying:
\begin{equation}
    \Phi(A)(x) \ge 0 \quad \text{a.e. } x \in \Omega, \quad \text{and} \quad \int_\Omega \Phi(A)(x) \dif x = \Tr(A) = 1.
\end{equation}
Each normalized matrix $A \in \mathcal{S}_+^n$ thus induces an absolutely continuous Borel probability measure on $\Omega$:
\begin{equation}
    \dif \mu_A(x) \coloneqq \Phi(A)(x) \dif x \in \mathcal{P}_{\mathrm{ac}}(\Omega).
\end{equation}
\end{definition}

A canonical family of such realizations is given by continuous kernel smoothing with a symmetric partition of unity $\{\psi_i\}_{i=1}^n \subset C^\infty(\Omega)$:
\begin{equation}
    \Phi_{\mathrm{kernel}}(A)(x) \coloneqq \sum_{i,j=1}^n A_{ij} \psi_i(x) \psi_j(x) = \boldsymbol{\psi}(x)^T A \boldsymbol{\psi}(x),
\end{equation}
where $\sum_{i=1}^n \int_\Omega \psi_i(x)^2 \dif x = 1$ and $\int_\Omega \psi_i(x)\psi_j(x)\dif x = \delta_{ij}$. Because $A \ge 0$, $\boldsymbol{\psi}(x)^T A \boldsymbol{\psi}(x) \ge 0$ for every $x \in \Omega$.

\subsection{Dimension-Free Matrix Wasserstein Distance}
Let $A \in \mathcal{S}_+^n$ and $B \in \mathcal{S}_+^m$ be two matrices of potentially \emph{different} dimensions $n \ne m$. Classical matrix analysis cannot define $\|A - B\|$. In our framework, their functional realizations induce probability measures $\mu_A, \mu_B \in \mathcal{P}(\Omega)$.

\begin{definition}[Functional $L^2$-Wasserstein Matrix Distance]
The functional quadratic Wasserstein distance between $A \in \mathcal{S}_+^n$ and $B \in \mathcal{S}_+^m$ is defined by:
\begin{equation}
    \mathcal{W}_2(A, B) \coloneqq \left( \inf_{\pi \in \Pi(\mu_A, \mu_B)} \int_{\Omega \times \Omega} \|x - y\|^2 \dif \pi(x, y) \right)^{1/2},
    \label{eq:wasserstein_matrix}
\end{equation}
where $\Pi(\mu_A, \mu_B)$ is the set of all Borel probability measures on $\Omega \times \Omega$ having marginals $\mu_A$ and $\mu_B$.
\end{definition}

By the Benamou--Brenier theorem \cite{benamou2000computational}, $\mathcal{W}_2(A, B)$ admits a dynamic fluid-mechanical formulation:
\begin{equation}
    \mathcal{W}_2^2(A, B) = \inf_{(\rho_t, \mathbf{v}_t)} \left\{ \int_0^1 \int_\Omega \|\mathbf{v}_t(x)\|^2 \rho_t(x) \dif x \dif t \right\},
\end{equation}
subject to the continuity equation and boundary conditions:
\begin{equation}
    \partial_t \rho_t + \nabla \cdot (\rho_t \mathbf{v}_t) = 0, \quad \rho_0 = \Phi(A), \quad \rho_1 = \Phi(B).
\end{equation}

\begin{theorem}[Metric Properties on Matrix Varieties]
\label{thm:wasserstein_properties}
Let $\mathcal{M} = \bigcup_{n=1}^\infty \mathcal{S}_+^n$ be the graded space of all normalized PSD matrices of arbitrary finite dimension. Then:
\begin{enumerate}[label=(\alph*)]
    \item $(\mathcal{M}, \mathcal{W}_2)$ is a pseudometric space. If $\Phi$ is injective on each $\mathcal{S}_+^n$, it defines a genuine metric on equivalence classes of matrices under realization kernel gauge equivalence.
    \item The distance $\mathcal{W}_2$ is continuous with respect to the $L^1(\Omega)$ norm:
    \begin{equation}
        \mathcal{W}_2(A, B) \le \sqrt{\diam(\Omega)} \|\Phi(A) - \Phi(B)\|_{L^1(\Omega)}^{1/2}.
    \end{equation}
    \item If $A(t)$ evolves under a continuous projected gradient flow $\dot{A} = -\nabla E(A)$ on $\mathcal{S}_+^n$, the curve $t \mapsto \Phi(A(t))$ is an absolutely continuous curve in the Wasserstein space $(\mathcal{P}_2(\Omega), \mathcal{W}_2)$ with metric velocity:
    \begin{equation}
        |\dot{\mu}_A|(t) \coloneqq \lim_{h \to 0} \frac{\mathcal{W}_2(A(t+h), A(t))}{|h|} \le C \|\dot{A}(t)\|_F.
    \end{equation}
\end{enumerate}
\end{theorem}

\begin{proof}
(a) Symmetry $\mathcal{W}_2(A, B) = \mathcal{W}_2(B, A)$ and non-negativity are immediate from \eqref{eq:wasserstein_matrix}. The triangle inequality follows from the gluing lemma for optimal couplings (Villani \cite{villani2009optimal}). 

(b) Let $\pi_0$ be the coupling with independent marginals plus diagonal concentration. By standard transportation inequalities:
\begin{equation}
    \mathcal{W}_2^2(\mu_A, \mu_B) \le \diam(\Omega)^2 \|\mu_A - \mu_B\|_{\mathrm{TV}} = \diam(\Omega)^2 \frac{1}{2} \|\Phi(A) - \Phi(B)\|_{L^1(\Omega)},
\end{equation}
from which the bound directly follows.

(c) Because $\Phi$ is smooth with bounded derivatives, $\left\|\frac{\dif}{\dif t}\Phi(A(t))\right\|_{L^1} \le C_1 \|\dot{A}(t)\|_F$. By the dynamic formulation, the continuity equation $\partial_t \rho + \nabla \cdot (\rho \mathbf{v}) = 0$ is solved by elliptic reconstruction with $\|\mathbf{v}\|_{L^2(\rho)} \le C_2 \|\partial_t \rho\|_{H^{-1}} \le C_3 \|\dot{A}\|_F$. Hence the metric derivative $|\dot{\mu}_A|(t) \le C \|\dot{A}(t)\|_F$.
\end{proof}

\subsection{Bakry--\'Emery Metric Curvature of a Matrix Realization}
Given a strictly positive functional realization $\Phi(A) \in C^2(\Omega)$ with $\Phi(A)(x) > 0$, we view $(\Omega, g, \dif \mu_A)$ as a weighted Riemannian metric measure space (smooth metric measure space) with background Euclidean metric $g_{ij} = \delta_{ij}$ and smooth potential:
\begin{equation}
    V_A(x) \coloneqq -\log \Phi(A)(x) \implies \dif \mu_A(x) = e^{-V_A(x)} \dif x.
\end{equation}
The weighted Laplacian (drift Laplacian / Witten Laplacian) acting on $f \in C^2(\Omega)$ is:
\begin{equation}
    \Delta_A f \coloneqq \Delta f - \langle \nabla V_A, \nabla f \rangle = \Delta f + \frac{\langle \nabla \Phi(A), \nabla f \rangle}{\Phi(A)}.
\end{equation}

\begin{definition}[Bakry--\'Emery Curvature Tensor and Invariant]
The infinite-dimensional Bakry--\'Emery Ricci tensor associated with the matrix $A$ is:
\begin{equation}
    \Ric_\infty(\Phi(A)) \coloneqq \Ric_g + \nabla^2 V_A = -\nabla^2 \log \Phi(A) = -\frac{\nabla^2 \Phi(A)}{\Phi(A)} + \frac{\nabla \Phi(A) \otimes \nabla \Phi(A)}{\Phi(A)^2}.
\end{equation}
The \emph{Bakry--\'Emery Curvature Invariant} $\kappa_{\mathrm{BE}}(A)$ of the matrix is the infimum of its spectrum:
\begin{equation}
    \kappa_{\mathrm{BE}}(A) \coloneqq \inf_{x \in \Omega, \|\xi\|=1} \left\langle \left( -\nabla^2 \log \Phi(A)(x) \right) \xi, \xi \right\rangle.
\end{equation}
\end{definition}

\begin{theorem}[Displacement Convexity and Log-Sobolev Inequality]
\label{thm:bakry_emery}
Let $A \in \mathcal{S}_+^n$ have smooth positive realization with $\kappa_{\mathrm{BE}}(A) \ge K > 0$. Then:
\begin{enumerate}[label=(\alph*)]
    \item \textbf{Displacement Convexity:} The Boltzmann--Shannon relative entropy functional $\mathrm{Ent}( \mu \mid \mu_A ) \coloneqq \int_\Omega \rho \log(\rho / \Phi(A)) \dif x$ is $K$-geodesically convex in the $L^2$-Wasserstein space $(\mathcal{P}_2(\Omega), \mathcal{W}_2)$.
    \item \textbf{Logarithmic Sobolev Inequality (LSI):} For every smooth function $f \in C^\infty(\Omega)$:
    \begin{equation}
        \int_\Omega f^2 \log\left(\frac{f^2}{\int_\Omega f^2 \dif\mu_A}\right) \dif\mu_A \le \frac{2}{K} \int_\Omega \|\nabla f\|^2 \dif\mu_A.
    \end{equation}
    \item \textbf{Exponential Convergence of Diffusion:} The heat flow $u_t = e^{t \Delta_A} u_0$ associated with the matrix $A$ contracts exponentially to the uniform state with rate $2K$:
    \begin{equation}
        \mathrm{Var}_{\mu_A}(e^{t \Delta_A} u_0) \le e^{-2Kt} \mathrm{Var}_{\mu_A}(u_0).
    \end{equation}
\end{enumerate}
\end{theorem}

\begin{proof}
By the Bakry--\'Emery criterion \cite{bakry1985diffusions}, the $\Gamma_2$ operator associated with the generator $\Delta_A$ satisfies:
\begin{equation}
    \Gamma_2(f, f) \coloneqq \frac{1}{2}\Delta_A \|\nabla f\|^2 - \langle \nabla f, \nabla \Delta_A f \rangle = \|\nabla^2 f\|_{\mathrm{HS}}^2 + \Ric_\infty(\Phi(A))(\nabla f, \nabla f).
\end{equation}
Since $\Ric_\infty(\Phi(A)) \ge K \cdot \mathrm{Id}$ everywhere, we have $\Gamma_2(f, f) \ge K \|\nabla f\|^2 = K \Gamma(f, f)$. By the classic Bakry--\'Emery theorem and Sturm--von Renesse \cite{von2005transport}, the $\Gamma_2 \ge K \Gamma$ condition is equivalent to the $K$-displacement convexity of the entropy along $\mathcal{W}_2$ geodesics, which immediately implies the Logarithmic Sobolev Inequality with constant $C_{\mathrm{LSI}} = 2/K$ and spectral gap $\lambda_1(-\Delta_A) \ge K$, completing the proof.
\end{proof}

% =========================================================================
\section{Persistent Homology, Topological Landscape, and Persistent Entropy}
\label{sec:persistent_homology}
% =========================================================================

\subsection{The Super-Level Set Filtration}
Let $A \in \R^{n \times n}$ be an arbitrary real matrix, and let $\Phi(A) \in C^0(\Omega)$ be its continuous functional realization on $\Omega = [0, 1]^d$. Let $M_A \coloneqq \max_{x \in \Omega} \Phi(A)(x)$ and $m_A \coloneqq \min_{x \in \Omega} \Phi(A)(x)$.

\begin{definition}[Continuous Matrix Filtration]
For each threshold $t \in [m_A, M_A]$, the continuous super-level set of the matrix $A$ is:
\begin{equation}
    X_t(A) \coloneqq \left\{ x \in \Omega \;\middle|\; \Phi(A)(x) \ge t \right\}.
\end{equation}
For $s \ge t$, we have the canonical inclusion $i_{s,t} \colon X_s(A) \hookrightarrow X_t(A)$, inducing a functorial persistence module over $(\R, \le)$:
\begin{equation}
    \mathcal{H}_k(A) \colon t \longmapsto H_k(X_t(A); \mathbb{F}),
\end{equation}
where $\mathbb{F}$ is a field (e.g., $\mathbb{Z}_2$ or $\R$) and $H_k$ denotes singular homology in degree $k \in \{0, 1, \dots, d-1\}$.
\end{definition}

By the Structure Theorem for persistence modules of tame functions (Edelsbrunner and Harer \cite{edelsbrunner2010computational}), $\mathcal{H}_k(A)$ decomposes into a multiset of interval modules, uniquely characterized by the \emph{Persistence Diagram}:
\begin{equation}
    \mathrm{Dgm}_k(\Phi(A)) \coloneqq \left\{ (b_i, d_i) \in \R^2 \;\middle|\; b_i > d_i, \; i \in I_k \right\} \cup \Delta,
\end{equation}
where $b_i$ is the birth value (a local maximum where a homology cycle appears), $d_i$ is the death value (where the cycle merges into an older cycle or fills in), and $\Delta = \{(x, x) \mid x \in \R\}$ is the diagonal with infinite multiplicity.

\subsection{Bottleneck Stability of Functional Tensors}
Let $\mathrm{Dgm}_k(f)$ and $\mathrm{Dgm}_k(g)$ be two persistence diagrams. The bottleneck distance is:
\begin{equation}
    d_B(\mathrm{Dgm}_k(f), \mathrm{Dgm}_k(g)) \coloneqq \inf_{\gamma \colon \mathrm{Dgm}_k(f) \to \mathrm{Dgm}_k(g)} \sup_{p \in \mathrm{Dgm}_k(f)} \|p - \gamma(p)\|_\infty,
\end{equation}
where $\gamma$ ranges over all bijections between the multisets.

\begin{theorem}[Bottleneck Stability for Functional Realizations]
\label{thm:bottleneck_stability}
Let $A, B \in \R^{n \times n}$ and let $\Phi$ be an $L^\infty$-bounded functional realization operator with Lipschitz constant $L_\Phi$:
\begin{equation}
    \|\Phi(A) - \Phi(B)\|_{L^\infty(\Omega)} \le L_\Phi \|A - B\|_F.
\end{equation}
Then, for every homology dimension $k \in \{0, 1, \dots, d-1\}$, the persistence diagrams satisfy:
\begin{equation}
    d_B(\mathrm{Dgm}_k(\Phi(A)), \mathrm{Dgm}_k(\Phi(B))) \le \|\Phi(A) - \Phi(B)\|_{L^\infty(\Omega)} \le L_\Phi \|A - B\|_F.
\end{equation}
\end{theorem}

\begin{proof}
Let $f = \Phi(A)$ and $g = \Phi(B)$. For any $\epsilon = \|f - g\|_{L^\infty(\Omega)}$, we have $f(x) \le g(x) + \epsilon$ and $g(x) \le f(x) + \epsilon$ for all $x \in \Omega$. 
This induces interlacing inclusions between the super-level sets:
\begin{equation}
    X_{t+\epsilon}(f) \subseteq X_t(g) \subseteq X_{t-\epsilon}(f).
\end{equation}
Applying the homology functor $H_k(\cdot; \mathbb{F})$ produces an algebraic $\epsilon$-interleaving between the persistence modules $\mathcal{H}_k(f)$ and $\mathcal{H}_k(g)$. By the Algebraic Stability Theorem of Chazal et al. \cite{chazal2016structure}, an $\epsilon$-interleaving implies that the bottleneck distance between the respective persistence diagrams is bounded by $\epsilon$:
\begin{equation}
    d_B(\mathrm{Dgm}_k(f), \mathrm{Dgm}_k(g)) \le \epsilon = \|f - g\|_{L^\infty(\Omega)}.
\end{equation}
Combining this with the Lipschitz continuity of $\Phi$ completes the proof.
\end{proof}

\subsection{The Persistent Entropy Invariant}
\begin{definition}[Persistent Entropy of a Matrix]
Let $\mathrm{Dgm}_k(\Phi(A)) = \{(b_i, d_i)\}_{i=1}^{N_k}$ be the off-diagonal points of the $k$-th persistence diagram. The lifespan of the $i$-th topological feature is $\ell_i \coloneqq b_i - d_i > 0$. The total persistent lifetime is $L_k(A) \coloneqq \sum_{i=1}^{N_k} \ell_i$. The persistent probability weights are $p_i \coloneqq \ell_i / L_k(A)$.
The \emph{$k$-th Persistent Entropy} of the matrix is:
\begin{equation}
    E_{\mathrm{pers}}^{(k)}(A) \coloneqq -\sum_{i=1}^{N_k} p_i \log p_i \ge 0.
\end{equation}
\end{definition}

\begin{theorem}[Isospectral Separation Theorem]
\label{thm:isospectral_separation}
There exist families of isospectral symmetric matrices $A, B \in \mathcal{S}^n$ such that:
\begin{equation}
    \sigma(A) = \sigma(B), \quad \Tr(A^p) = \Tr(B^p) \; \forall p \in \mathbb{N}, \quad \|A\|_F = \|B\|_F,
\end{equation}
and yet their functional topological invariants strictly differ:
\begin{equation}
    \mathrm{Dgm}_1(\Phi(A)) \ne \mathrm{Dgm}_1(\Phi(B)) \quad \text{and} \quad E_{\mathrm{pers}}^{(1)}(A) \ne E_{\mathrm{pers}}^{(1)}(B).
\end{equation}
\end{theorem}

\begin{proof}
We construct an explicit counter-example. Let $n = 4$. Consider the continuous step-realization on $\Omega = [0, 1]^2$ partitioned into a $4 \times 4$ grid of cells $\Omega_{ij} = [\frac{i-1}{4}, \frac{i}{4}] \times [\frac{j-1}{4}, \frac{j}{4}]$.
Let $A$ be a block-diagonal matrix:
\begin{equation}
    A = \begin{pmatrix} 2 & 1 & 0 & 0 \\ 1 & 2 & 0 & 0 \\ 0 & 0 & -1 & 0 \\ 0 & 0 & 0 & -1 \end{pmatrix},
\end{equation}
and let $P$ be the permutation matrix transposing rows/columns 2 and 3:
\begin{equation}
    B = P A P^T = \begin{pmatrix} 2 & 0 & 1 & 0 \\ 0 & -1 & 0 & 0 \\ 1 & 0 & 2 & 0 \\ 0 & 0 & 0 & -1 \end{pmatrix}.
\end{equation}
Since $B = P A P^T$, $\sigma(A) = \sigma(B) = \{3, 1, -1, -1\}$ and all algebraic polynomial traces are identical.
However, in the spatial domain $\Omega = [0, 1]^2$, the functional realization $\Phi(A)(x, y)$ concentrates high values ($2$ and $1$) in the contiguous $2 \times 2$ top-left subregion $\Omega_{1..2, 1..2}$, forming a single simply-connected spatial peak whose super-level sets $X_t(A)$ for $t \in (0, 2]$ have the homotopy type of a point ($\beta_0 = 1, \beta_1 = 0$).

Conversely, in $B$, the entries of magnitude $\ge 1$ are located at $(1,1), (3,3), (1,3), (3,1)$, which form four isolated corners separated by negative valleys (entries $-1$ at $(2,2)$ and $(4,4)$ and zeros at other cells). For $t \in (0, 1.5]$, the super-level set $X_t(B)$ forms a non-contractible topological ring enclosing the central low-value pit at cell $(2,2)$, creating a 1-dimensional homology cycle:
\begin{equation}
    H_1(X_t(B); \mathbb{Z}_2) \cong \mathbb{Z}_2 \implies \beta_1(X_t(B)) = 1,
\end{equation}
whereas $H_1(X_t(A); \mathbb{Z}_2) = 0$ for all $t$.
Hence, $\mathrm{Dgm}_1(\Phi(B))$ contains a non-trivial point $(b, d) \approx (2, 0)$ of persistence length $\ell = 2$, yielding $E_{\mathrm{pers}}^{(1)}(B) > 0$, while $\mathrm{Dgm}_1(\Phi(A)) = \emptyset$ and $E_{\mathrm{pers}}^{(1)}(A) = 0$. This rigorously proves strict topological separation.
\end{proof}

% =========================================================================
\section{Geometric Measure Theory: Willmore Bending and Isoperimetric Deficits}
\label{sec:willmore_gmt}
% =========================================================================

\subsection{Level-Set Varifolds and Mean Curvature Fields}
Let $\Phi(A) \in C^2(\Omega)$ be a smooth realization on a compact Riemannian domain $(\Omega, g)$. For almost every $t \in \R$, Sard's theorem guarantees that the level set
\begin{equation}
    \Sigma_t \coloneqq \left\{ x \in \Omega \;\middle|\; \Phi(A)(x) = t \right\}
\end{equation}
is a smooth $(d-1)$-dimensional embedded submanifold of $\Omega$. The unit normal vector field pointing towards increasing values is:
\begin{equation}
    \mathbf{n}(x) \coloneqq \frac{\nabla \Phi(A)(x)}{\|\nabla \Phi(A)(x)\|}, \quad x \in \Sigma_t.
\end{equation}

\begin{definition}[Mean Curvature Field of a Functional Realization]
The scalar mean curvature $H_t(x)$ of the level hypersurface $\Sigma_t$ at $x$ is given by the tangential divergence of the normal field:
\begin{equation}
    H_t(x) \coloneqq -\operatorname{div}\left( \frac{\nabla \Phi(A)(x)}{\|\nabla \Phi(A)(x)\|} \right) = -\frac{\Delta \Phi(A)}{\|\nabla \Phi(A)\|} + \frac{\Hess \Phi(A)(\nabla \Phi(A), \nabla \Phi(A))}{\|\nabla \Phi(A)\|^3}.
\end{equation}
\end{definition}

\subsection{The Integrated Willmore Bending Energy}
In classical differential geometry, the Willmore functional of a 2D surface $\Sigma \subset \R^3$ is $\int_\Sigma H^2 \dif A$, measuring its total elastic bending energy. We extend this to matrix realizations by integrating over all level sets.

\begin{definition}[Integrated Willmore Energy of a Matrix Realization]
The \emph{Integrated Willmore Bending Energy} $\mathcal{W}(\Phi(A))$ of a matrix $A$ is defined by:
\begin{equation}
    \mathcal{W}(\Phi(A)) \coloneqq \int_{-\infty}^\infty \left( \int_{\Sigma_t \cap \Omega} H_t(x)^2 \dif \mathcal{H}^{d-1}(x) \right) \dif t.
\end{equation}
Applying the Coarea formula ($\dif t \dif \mathcal{H}^{d-1} = \|\nabla \Phi(A)\| \dif x$), $\mathcal{W}$ admits the explicit volume integral representation:
\begin{equation}
    \mathcal{W}(\Phi(A)) = \int_\Omega \left| \operatorname{div}\left(\frac{\nabla \Phi(A)}{\|\nabla \Phi(A)\|}\right) \right|^2 \|\nabla \Phi(A)\| \dif x.
    \label{eq:willmore_volume}
\end{equation}
\end{definition}

\begin{theorem}[Conformal Invariance and Regularity Bounds]
\label{thm:willmore_invariance}
Let $d = 3$ and $\Omega \subset \R^3$. 
\begin{enumerate}[label=(\alph*)]
    \item \textbf{Conformal Invariance of Level Energy:} If $\psi \colon \R^3 \to \R^3$ is a conformal diffeomorphism (a composition of translations, dilations, rotations, and sphere inversions), then for each level set $\Sigma_t$, the trace-free curvature Willmore integrand is conformally invariant:
    \begin{equation}
        \int_{\psi(\Sigma_t)} (H^2 - K) \dif \mathcal{H}^2 = \int_{\Sigma_t} (H^2 - K) \dif \mathcal{H}^2,
    \end{equation}
    where $K$ is the Gaussian curvature.
    \item \textbf{High-Frequency Penalization:} If $A_k$ is a sequence of matrices representing high-frequency checkerboard oscillations of frequency $k$ and fixed amplitude $\epsilon$, their Frobenius norm remains bounded ($\|A_k\|_F = O(1)$), but their Willmore energy diverges quadratically:
    \begin{equation}
        \mathcal{W}(\Phi(A_k)) = \Theta(k^2) \longrightarrow \infty \quad \text{as } k \to \infty.
    \end{equation}
\end{enumerate}
\end{theorem}

\begin{proof}
(a) Follows directly from the classical theorem of Willmore and Chen \cite{willmore1993riemannian} on the conformal invariance of $\int (H^2 - K) \dif A$ for 2-surfaces in $\R^3$. 

(b) For a checkerboard realization of wavelength $\lambda \sim 1/k$, consider:
\begin{equation}
    \Phi(A_k)(x) \sim \epsilon \sin(k x_1)\sin(k x_2).
\end{equation}
Then $\|\nabla \Phi\| \sim \epsilon k$, and the Hessian terms scale as $\Hess \Phi \sim \epsilon k^2$. 
The mean curvature scales as:
\begin{equation}
    H_t(x) \sim \frac{\epsilon k^2}{\epsilon k} = k.
\end{equation}
Substituting into the volume representation \eqref{eq:willmore_volume}:
\begin{equation}
    \mathcal{W}(\Phi(A_k)) \sim \int_\Omega (k)^2 (\epsilon k) \dif x \sim \epsilon k^3.
\end{equation}
With normalized amplitude $\epsilon \sim 1/k$ to keep gradient $L^1$ bounded, $\mathcal{W} \sim k^2 \to \infty$. Thus, Willmore energy acts as a strong second-order regularizer that penalizes geometric wrinkling far more strictly than the $H^1$ Dirichlet energy.
\end{proof}

\subsection{The Matrix Isoperimetric Profile}
\begin{definition}[Isoperimetric Profile]
For a positive functional realization $\Phi(A) \ge 0$ inducing the measure $\mu_A$, the \emph{Isoperimetric Profile} $\mathcal{I}_A \colon (0, 1) \to \R_+$ is:
\begin{equation}
    \mathcal{I}_A(v) \coloneqq \inf \left\{ \Per_{\mu_A}(E) \;\middle|\; E \subset \Omega \text{ Borel}, \; \mu_A(E) = v \right\},
\end{equation}
where $\Per_{\mu_A}(E) \coloneqq \int_{\partial^* E} \Phi(A)(x) \dif \mathcal{H}^{d-1}(x)$ is the weighted perimeter.
The \emph{Cheeger Constant} of the matrix is:
\begin{equation}
    h(A) \coloneqq \inf_{0 < \mu_A(E) \le 1/2} \frac{\Per_{\mu_A}(E)}{\mu_A(E)}.
\end{equation}
\end{definition}
By the Cheeger--Federer inequality, the first non-trivial eigenvalue of the weighted drift Laplacian satisfies $\lambda_1(-\Delta_A) \ge \frac{1}{4} h(A)^2$. Thus, $h(A)$ provides a direct geometric measure of the **spectral bottleneck** in the matrix graph.

% =========================================================================
\section{Micro-Local Phase-Space Analysis and Fractional Besov Regularity}
\label{sec:microlocal}
% =========================================================================

\subsection{Continuous Gabor Wavepacket Transform}
To localize not merely the spatial position $x \in \Omega$ of matrix features, but their frequency and direction $\xi \in \R^d \setminus \{0\}$, we lift the functional realization $\Phi(A)$ to the cotangent bundle $T^*\Omega \cong \Omega \times \R^d$.

Let $g \in C_c^\infty(\R^d)$ be a symmetric Gaussian window function with $\|g\|_{L^2} = 1$.
\begin{definition}[Continuous Matrix Wavepacket Transform]
For $x \in \Omega$ and $\xi \in \R^d$, the Gabor transform of the realization $\Phi(A) \in L^2(\Omega)$ is:
\begin{equation}
    \mathcal{V}_{\Phi(A)}(x, \xi) \coloneqq \int_{\R^d} \Phi(A)(y) g(y - x) e^{-i \langle \xi, y \rangle} \dif y.
\end{equation}
The associated \emph{Phase-Space Energy Density} (Wigner--Radon measure) is:
\begin{equation}
    \mathcal{E}_A(x, \xi) \coloneqq \left| \mathcal{V}_{\Phi(A)}(x, \xi) \right|^2 \ge 0.
\end{equation}
\end{definition}

\subsection{The Matrix Wavefront Set}
In micro-local analysis (Hörmander \cite{hormander2003analysis}), the wavefront set characterizes the directions in which a distribution fails to be smooth.

\begin{definition}[Matrix Wavefront Set]
A point $(x_0, \xi_0) \in \Omega \times (\R^d \setminus \{0\})$ does \emph{not} belong to the wavefront set $\WF(\Phi(A))$ if there exist a neighborhood $U$ of $x_0$ and an open cone $\Gamma \subset \R^d \setminus \{0\}$ containing $\xi_0$ such that for every $N \in \mathbb{N}$:
\begin{equation}
    \sup_{x \in U, \xi \in \Gamma} (1 + \|\xi\|)^N \left| \mathcal{V}_{\Phi(A)}(x, \xi) \right| < \infty.
\end{equation}
The \emph{Matrix Wavefront Set} is the complement:
\begin{equation}
    \WF(\Phi(A)) \subset T^*\Omega \setminus \{0\}.
\end{equation}
\end{definition}

\begin{theorem}[Interface and Shock Direction Detection]
\label{thm:wavefront_interface}
Let $A \in \R^{n \times n}$ be a piecewise constant step realization with jump discontinuity across a smooth interface $\Gamma \subset \Omega$ with unit normal $\mathbf{n}_\Gamma(x)$. Then:
\begin{equation}
    \WF(\Phi(A)) = N^*(\Gamma) \coloneqq \left\{ (x, \xi) \in \Gamma \times (\R^d \setminus \{0\}) \;\middle|\; \xi = \lambda \mathbf{n}_\Gamma(x), \; \lambda \ne 0 \right\},
\end{equation}
where $N^*(\Gamma)$ is the conormal bundle of the interface $\Gamma$.
\end{theorem}

\begin{proof}
Away from $\Gamma$, $\Phi(A)$ is locally constant, hence smooth ($C^\infty$). By standard integration by parts against the Gaussian window $g(y - x)$, $\mathcal{V}_{\Phi(A)}(x, \xi)$ decays faster than any polynomial power $(1 + \|\xi\|)^{-N}$ for any $x \notin \Gamma$. Hence $\WF(\Phi(A)) \subseteq \Gamma \times (\R^d \setminus \{0\})$.

At $x_0 \in \Gamma$, choose local coordinates such that $\Gamma = \{y_1 = 0\}$ and $\mathbf{n}_\Gamma(x_0) = (1, 0, \dots, 0)$. The jump in $\Phi(A)$ is $c \, H(y_1)$ where $H$ is the Heaviside step function and $c \ne 0$. The partial Fourier transform in the tangential directions $(y_2, \dots, y_d)$ decays rapidly because the jump profile is invariant along the tangent plane $T_{x_0}\Gamma$. In the normal direction $\xi_1$, the 1D Fourier transform of the Heaviside jump is:
\begin{equation}
    \widehat{H}(\xi_1) = \frac{1}{i \xi_1} + \pi \delta(\xi_1).
\end{equation}
This decays only as $O(\|\xi_1\|^{-1})$, violating super-algebraic decay for all $\xi_1 \ne 0$. Therefore, $(x_0, \xi) \in \WF(\Phi(A))$ if and only if $\xi \parallel \mathbf{n}_\Gamma(x_0)$, which precisely defines the conormal bundle $N^*(\Gamma)$.
\end{proof}

\subsection{Fractional Besov Exponents}
To quantify the exact fractional smoothness of singular tensors, we use the Gagliardo--Besov seminorm:
\begin{equation}
    [\Phi(A)]_{B^s_{p,\infty}} \coloneqq \sup_{h \in \R^d, 0 < \|h\| \le 1} \frac{\|\Phi(A)(\cdot + h) - \Phi(A)(\cdot)\|_{L^p(\Omega_h)}}{\|h\|^s}, \quad s \in (0, 1).
\end{equation}

\begin{definition}[Critical Fractional Regularity Index]
The \emph{Critical Regularity Index} $s^*(A)$ of the matrix realization is:
\begin{equation}
    s^*(A) \coloneqq \sup \left\{ s \in (0, 1) \;\middle|\; [\Phi(A)]_{B^s_{1,\infty}} < \infty \right\}.
\end{equation}
\end{definition}
For smooth realizations, $s^*(A) = 1$. For piecewise constant step matrices, $s^*(A) = 1/p$ (in $L^p$) and $s^*(A) = 1$ in $B^1_{1,\infty} = \BV$. For fractal or self-similar matrices (such as Pascal or Sierpi\'nski matrices), $s^*(A) = d - d_{\mathrm{fractal}}$, providing a continuous algebraic detector of fractal dimension.

% =========================================================================
\section{Non-Commutative Spectral Invariants and Quantum Information Geometry}
\label{sec:noncommutative}
% =========================================================================

\subsection{The Canonical Dirac Spectral Triple}
We lift the functional realization framework to Alain Connes' Non-Commutative Geometry \cite{connes1994noncommutative}.
Let $(\Omega, g)$ be a compact $d$-dimensional spin Riemannian manifold. We define the canonical spectral triple $(\mathcal{A}_\Phi, \mathcal{H}, \mathcal{D})$:
\begin{itemize}
    \item $\mathcal{H} \coloneqq L^2(\Omega, \mathbb{S})$: the Hilbert space of square-integrable Dirac spinors on $\Omega$.
    \item $\mathcal{D} \coloneqq i \gamma^\mu \nabla_\mu^\mathbb{S}$: the classical Atiyah--Singer Dirac operator, which is self-adjoint with discrete spectrum $|\lambda_n| \to \infty$.
    \item $\mathcal{A}_\Phi$: the $*$-algebra generated by multiplication operators $M_{\Phi(A)} \colon \psi \mapsto \Phi(A) \cdot \psi$ for all smooth matrix realizations $A \in \R^{n \times n}$.
\end{itemize}

\subsection{The Connes--Dixmier Conformal Volume Trace}
On $\mathcal{H}$, the operator $|\mathcal{D}|^{-d}$ is in the Dixmier ideal $\mathcal{L}^{1,\infty}(\mathcal{H})$.

\begin{definition}[Non-Commutative Matrix Trace]
For $A \in \R^{n \times n}$, the \emph{Connes--Dixmier Matrix Invariant} is defined by:
\begin{equation}
    \Tr_\omega \left( \Phi(A) |\mathcal{D}|^{-d} \right) \coloneqq \lim_{N \to \infty} \frac{1}{\log N} \sum_{n=1}^N \mu_n \left( \Phi(A) |\mathcal{D}|^{-d} \right),
\end{equation}
where $\mu_n(\cdot)$ are the singular values of the compact operator.
\end{definition}

\begin{theorem}[Non-Commutative Integration Formula]
\label{thm:dixmier_integration}
Let $\Phi(A) \in C^\infty(\Omega)$. The non-commutative trace extracts the continuous geometric integral of the realization:
\begin{equation}
    \Tr_\omega \left( \Phi(A) |\mathcal{D}|^{-d} \right) = \frac{2^{\lfloor d/2 \rfloor} \Omega_d}{d (2\pi)^d} \int_\Omega \Phi(A)(x) \dif \mathrm{vol}_g(x),
\end{equation}
where $\Omega_d = \frac{2 \pi^{d/2}}{\Gamma(d/2)}$ is the volume of the $(d-1)$-sphere $\mathbb{S}^{d-1}$.
\end{theorem}

\begin{proof}
This is an application of Connes' trace theorem \cite{connes1994noncommutative} to the order $-d$ pseudo-differential operator $T = \Phi(A) |\mathcal{D}|^{-d}$. 
The principal symbol of $\mathcal{D}$ is $\sigma_1(\mathcal{D})(x, \xi) = \gamma(\xi) = \sum_{j=1}^d \gamma^j \xi_j$. 
Since $\gamma(\xi)^2 = \|\xi\|^2 \Id$, the principal symbol of $|\mathcal{D}|^{-d}$ is $\|\xi\|^{-d} \Id_{\mathbb{S}}$. 
Thus, the principal symbol of $T$ is:
\begin{equation}
    \sigma_{-d}(T)(x, \xi) = \Phi(A)(x) \|\xi\|^{-d} \Id_{\mathbb{S}}.
\end{equation}
The Wodzicki residue (which coincides with the Dixmier trace up to normalization) is the integral of the trace of the symbol over the cosphere bundle $S^*\Omega = \{(x, \xi) \mid \|\xi\|_g = 1\}$:
\begin{equation}
    \Tr_\omega(T) = \frac{1}{d (2\pi)^d} \int_{S^*\Omega} \Tr_{\mathbb{S}} \left( \sigma_{-d}(T)(x, \xi) \right) \dif S(x, \xi).
\end{equation}
Since $\Tr_{\mathbb{S}}(\Id) = 2^{\lfloor d/2 \rfloor}$, the $\xi$-integration over the unit sphere yields $\Omega_d$, and the $x$-integration yields $\int_\Omega \Phi(A)(x) \dif \mathrm{vol}_g(x)$, proving the formula.
\end{proof}

\subsection{Quantized Cyclic Cocycles and Topological Indices}
The commutator of $\mathcal{D}$ with the functional realization is:
\begin{equation}
    [\mathcal{D}, \Phi(A)] = i \gamma^\mu (\partial_\mu \Phi(A)) = i \gamma(\dif \Phi(A)),
\end{equation}
which represents the Clifford-valued gradient. For a tuple of $d$ matrices $(A_0, A_1, \dots, A_d)$, we define the \emph{Non-Commutative Matrix Winding Invariant}:
\begin{equation}
    \tau_d(A_0, A_1, \dots, A_d) \coloneqq \Tr_\omega \left( \Phi(A_0) [\mathcal{D}, \Phi(A_1)] [\mathcal{D}, \Phi(A_2)] \cdots [\mathcal{D}, \Phi(A_d)] |\mathcal{D}|^{-d} \right).
\end{equation}
When the maps $\Phi(A_j)$ define coordinates of a map from $\Omega$ into a target sphere, $\tau_d$ computes the integer topological Brouwer degree (winding number), providing a **quantized topological charge** on the matrix manifold.

% =========================================================================
\section{Quantum Information Geometry of Continuous Tensor Networks}
\label{sec:quantum_info}
% =========================================================================

\subsection{Fiber-Bundle Tensor Networks}
In quantum physics, a continuous Matrix Product State (cMPS) \cite{verstraete2010continuous} on $\Omega$ assigns to each point $x \in \Omega$ an auxiliary bond Hilbert space $\mathbb{C}^\chi$ and local field operators $Q(x), R(x) \in \mathrm{End}(\mathbb{C}^\chi)$. 
This endows the functional realization with a canonical **fiber bundle structure**:
\begin{equation}
    \mathcal{E} \coloneqq \Omega \times \mathrm{End}(\mathbb{C}^\chi) \longrightarrow \Omega,
\end{equation}
where each fiber above $x$ contains the normalized auxiliary density matrix $\rho(x) \in \mathcal{S}_+(\mathbb{C}^\chi)$, $\Tr_\chi(\rho(x)) = 1$.

\subsection{The Integrated Quantum Fisher Information (QFI) Volume}
Let the state $\rho(x)$ vary smoothly across $\Omega$. In quantum estimation theory, the distance between neighboring states $\rho(x)$ and $\rho(x + \dif x)$ is governed by the Quantum Fisher Information metric:
\begin{equation}
    g^{\mathrm{QFI}}_{\mu\nu}(x) \coloneqq \frac{1}{2} \Tr_\chi \left( \rho(x) \{ L_\mu(x), L_\nu(x) \} \right),
\end{equation}
where $\{A, B\} = AB + BA$, and $L_\mu(x)$ is the Symmetric Logarithmic Derivative (SLD) defined implicitly by:
\begin{equation}
    \partial_\mu \rho(x) = \frac{1}{2} \left( \rho(x) L_\mu(x) + L_\mu(x) \rho(x) \right).
\end{equation}

\begin{definition}[Integrated QFI Volume]
The \emph{Integrated Quantum Fisher Information Volume} of a tensor network state $\mathcal{T}$ is:
\begin{equation}
    \mathrm{Vol}_{\mathrm{QFI}}(\mathcal{T}) \coloneqq \int_\Omega \sqrt{\det g^{\mathrm{QFI}}(x)} \dif^d x.
\end{equation}
\end{definition}

\begin{theorem}[Positivity and Entanglement Faithfulness]
\label{thm:qfi_volume}
Let $\mathcal{T}$ be a continuous matrix product state. Then:
\begin{enumerate}[label=(\alph*)]
    \item $\mathrm{Vol}_{\mathrm{QFI}}(\mathcal{T}) \ge 0$.
    \item $\mathrm{Vol}_{\mathrm{QFI}}(\mathcal{T}) = 0$ if and only if $\rho(x)$ is spatially constant across $\Omega$ modulo unitary gauge transformations $U(x) \in U(\chi)$ (i.e., the tensor network is a trivial, unentangled spatial product state).
    \item $\mathrm{Vol}_{\mathrm{QFI}}(\mathcal{T})$ is invariant under local spatial reparametrizations $\mathrm{Diff}^+(\Omega)$.
\end{enumerate}
\end{theorem}

\begin{proof}
(a) The SLD operator $L_\mu$ is Hermitian, and $\rho(x)$ is PSD. Thus, $g^{\mathrm{QFI}}(x)$ is a positive semi-definite matrix at every $x \in \Omega$, which implies $\det g^{\mathrm{QFI}}(x) \ge 0$, and hence $\mathrm{Vol}_{\mathrm{QFI}}(\mathcal{T}) \ge 0$.

(b) If $\mathrm{Vol}_{\mathrm{QFI}}(\mathcal{T}) = 0$, then $\det g^{\mathrm{QFI}}(x) = 0$ almost everywhere. By the fundamental property of the Bures metric \cite{bengtsson2017geometry}, $g^{\mathrm{QFI}}_{\mu\mu}(x) = 0$ if and only if $\partial_\mu \rho(x) = 0$. Hence $\rho(x)$ is spatially homogeneous, representing a tensor state with no spatial variation of entanglement Cauchy data.

(c) Under a change of coordinates $y = \phi(x)$, the metric tensor transforms as $g^{\mathrm{QFI}}_y = (J^{-1})^T g^{\mathrm{QFI}}_x J^{-1}$ where $J = \dif \phi / \dif x$. Thus $\sqrt{\det g^{\mathrm{QFI}}_y} \dif^d y = \sqrt{\det g^{\mathrm{QFI}}_x} \dif^d x$, proving diffeomorphism invariance.
\end{proof}

\subsection{Continuous Entanglement Contour and Centroid}
The scalar von Neumann entropy of the bond state is $S_{\mathrm{vN}}(x) = -\Tr_\chi(\rho(x) \log \rho(x))$. We define the geometric density of quantum entanglement:

\begin{definition}[Entanglement Contour and Centroid]
The \emph{Continuous Entanglement Contour} is the volume-weighted scalar density:
\begin{equation}
    \mathcal{S}_{\mathrm{cont}}(x) \coloneqq -\Tr_\chi(\rho(x) \log \rho(x)) \cdot \sqrt{\det g^{\mathrm{QFI}}(x)}.
\end{equation}
The \emph{Total Geometric Entanglement} is $S_{\mathrm{geom}}(\mathcal{T}) \coloneqq \int_\Omega \mathcal{S}_{\mathrm{cont}}(x) \dif^d x$.
If $S_{\mathrm{geom}}(\mathcal{T}) > 0$, the \emph{Entanglement Spatial Centroid} $\bar{x}_{\mathrm{ent}} \in \Omega$ and \emph{Spatial Entanglement Dispersion} $\Sigma_{\mathrm{ent}}^2$ are:
\begin{equation}
    \bar{x}_{\mathrm{ent}} \coloneqq \frac{1}{S_{\mathrm{geom}}(\mathcal{T})} \int_\Omega x \, \mathcal{S}_{\mathrm{cont}}(x) \dif^d x, \quad \Sigma_{\mathrm{ent}}^2 \coloneqq \frac{1}{S_{\mathrm{geom}}(\mathcal{T})} \int_\Omega \|x - \bar{x}_{\mathrm{ent}}\|^2 \, \mathcal{S}_{\mathrm{cont}}(x) \dif^d x.
\end{equation}
\end{definition}
These invariants localize **where** in physical space the entanglement of the tensor variety is concentrated, solving the spatial blind spot of finite-dimensional multilinear entanglement measures.

% =========================================================================
\section{Synthesis, Taxonomy of Invariants, and Applications}
\label{sec:synthesis}
% =========================================================================

\subsection{The Master Taxonomy of Functional Tensor Invariants}
Table \ref{tab:master_taxonomy} summarizes the mathematical taxonomy of all invariants developed in this treatise alongside their foundational counterparts from Volume I.

\begin{table}[h!]
\centering
\small
\renewcommand{\arraystretch}{1.25}
\begin{tabularx}{\textwidth}{@{} l p{2.8cm} p{3.2cm} X @{}}
\toprule
\textbf{Invariant} & \textbf{Mathematical Family} & \textbf{Formal Definition} & \textbf{Physical / Computational Role} \\
\midrule
$\mathcal{W}_2(A, B)$ & Optimal Transport & $\inf_\pi \int \|x-y\|^2 \dif \pi$ & Geodesic distance between matrices of unequal dimensions ($n \ne m$). \\
\midrule
$\kappa_{\mathrm{BE}}(A)$ & Metric Measure Geometry & $\inf_\xi \langle (-\nabla^2 \log \Phi) \xi, \xi \rangle$ & Ricci curvature, log-Sobolev constant, dissipation rate of flows. \\
\midrule
$\mathrm{Dgm}_k(A)$ & Topological Data Analysis & $\{(b_i, d_i) \mid H_k(X_t) \ne 0\}$ & Multi-scale topological persistence of cavities; breaks isospectrality. \\
\midrule
$E_{\mathrm{pers}}(A)$ & Persistent Entropy & $-\sum p_i \log p_i$ & Quantifies topological disorder and complexity in entry landscapes. \\
\midrule
$\mathcal{W}(\Phi(A))$ & Geometric Measure Theory & $\int \int_{\Sigma_t} H_t^2 \dif \mathcal{H}^{d-1} \dif t$ & Total Willmore bending; penalizes second-order surface wrinkling. \\
\midrule
$\mathcal{I}_A(v)$ & GMT / Isoperimetry & $\inf \{ \Per(E) \mid \mu_A(E) = v \}$ & Cheeger constant, matrix bottlenecks, conductance bounds. \\
\midrule
$\WF(\Phi(A))$ & Micro-Local Analysis & $\Sigma_{\mathrm{sing}} \subset T^*\Omega \setminus \{0\}$ & Exact spatial location and directional orientation of shock interfaces. \\
\midrule
$s^*(A)$ & Fractional Analysis & $\sup \{ s \mid [\Phi]_{B^s_{1,\infty}} < \infty \}$ & Critical Besov exponent; detects fractal dimension of entries. \\
\midrule
$\Tr_\omega$ & Non-Commutative Geom. & $\Tr_\omega(\Phi |\mathcal{D}|^{-d})$ & Conformal volume integral via Atiyah--Singer Dirac spectrum. \\
\midrule
$\mathrm{Vol}_{\mathrm{QFI}}(\mathcal{T})$ & Quantum Information & $\int_\Omega \sqrt{\det g^{\mathrm{QFI}}(x)} \dif^d x$ & Total distinguishability and parameter sensitivity of tensor states. \\
\midrule
$\bar{x}_{\mathrm{ent}}, \Sigma_{\mathrm{ent}}^2$ & Entanglement Contour & $\int x \mathcal{S}_{\mathrm{cont}}(x) \dif x$ & Physical center of mass and spatial spread of quantum entanglement. \\
\bottomrule
\end{tabularx}
\caption{Comprehensive taxonomy of advanced geometric and topological invariants on functional realizations of matrices and tensors.}
\label{tab:master_taxonomy}
\end{table}

\subsection{Key Scientific Applications}
\begin{enumerate}[label=(\alph*),leftmargin=*]
    \item \textbf{Transfer Learning and Dimension-Free Neural Weight Geometry:} In deep neural networks, comparing weight matrices $W_1 \in \R^{d_{\mathrm{in}} \times d_{\mathrm{out}}}$ and $W_2 \in \R^{d'_{\mathrm{in}} \times d'_{\mathrm{out}}}$ of different architectures has historically required ad-hoc projection or padding. The functional Wasserstein distance $\mathcal{W}_2(\Phi(W_1), \Phi(W_2))$ provides an intrinsic, continuous distance metric on network representations, enabling true architectural transfer metrics.
    \item \textbf{Topological Regularization of Deep Learning Landscapes:} Traditional weight decay penalizes $\|W\|_F^2 = \sum W_{ij}^2$, which is blind to spatial correlations. Adding the persistent entropy $E_{\mathrm{pers}}(\Phi(W))$ or Willmore energy $\mathcal{W}(\Phi(W))$ as a loss regularizer forces weight matrices to maintain smooth, low-curvature level sets, preventing high-frequency adversarial vulnerability.
    \item \textbf{Quantum Phase Transitions in Continuous Tensor Networks:} In continuous matrix product states, quantum critical points are signaled by the divergence of correlation lengths. While the scalar von Neumann entropy $S_{\mathrm{vN}}$ often exhibits logarithmic scaling, the integrated QFI volume $\mathrm{Vol}_{\mathrm{QFI}}(\mathcal{T})$ and the spatial entanglement dispersion $\Sigma_{\mathrm{ent}}^2$ detect the precise spatial spreading of quantum fluctuations at criticality.
\end{enumerate}

% =========================================================================
\section{Conclusion and Future Horizons}
\label{sec:conclusion}
% =========================================================================

In this work, we have demonstrated that embedding discrete matrices and tensors into continuous function spaces via Functional Realization Operators is not merely a technical interpolation, but a profound paradigm shift in multilinear algebra. By mobilizing optimal transport, persistent homology, geometric measure theory, micro-local analysis, non-commutative geometry, and quantum information geometry, we have uncovered a vast universe of invariants that are completely invisible to classical spectrum and matrix rank.

These invariants resolve fundamental open problems in dimensional incommensurability, isospectral ambiguity, high-frequency wrinkling, and spatial entanglement localization. Future research will focus on the complete machine-checked formalization of these theorems within interactive proof assistants such as \textsc{Lean 4}, expanding the open verification corpus and establishing a new rigorous foundation for computational and physical multilinear geometry.

% =========================================================================
\begin{thebibliography}{99}
% =========================================================================

\bibitem{bakry1985diffusions}
D.~Bakry and M.~{\'E}mery,
\emph{Diffusions hypercontractives},
S{\'e}minaire de Probabilit{\'e}s XIX 1983/84, Lecture Notes in Math., vol.~1123, Springer, Berlin, 1985, pp.~177--206.

\bibitem{benamou2000computational}
J.-D.~Benamou and Y.~Brenier,
\emph{A computational fluid mechanics solution to the Monge-Kantorovich mass transfer problem},
Numer. Math. \textbf{84} (2000), no.~3, 375--393.

\bibitem{bengtsson2017geometry}
I.~Bengtsson and K.~{\.Z}yczkowski,
\emph{Geometry of Quantum States: An Introduction to Quantum Entanglement},
2nd ed., Cambridge University Press, Cambridge, 2017.

\bibitem{chazal2016structure}
F.~Chazal, V.~de Silva, M.~Glisse, and S.~Oudot,
\emph{The Structure and Stability of Persistence Modules},
SpringerBriefs in Mathematics, Springer, Cham, 2016.

\bibitem{connes1994noncommutative}
A.~Connes,
\emph{Noncommutative Geometry},
Academic Press, San Diego, CA, 1994.

\bibitem{edelsbrunner2010computational}
H.~Edelsbrunner and J.~L.~Harer,
\emph{Computational Topology: An Introduction},
American Mathematical Society, Providence, RI, 2010.

\bibitem{hormander2003analysis}
L.~H{\"o}rmander,
\emph{The Analysis of Linear Partial Differential Operators I: Distribution Theory and Fourier Analysis},
Classics in Mathematics, Springer-Verlag, Berlin, 2003.

\bibitem{silvafilho2026beyond}
R.~M.~Silva-Filho,
\emph{Beyond the Spectrum: Functional Realizations of Matrices and Tensors, Emerging Invariants, and Geometric Measures},
Treatise on Multilinear Geometry and Quantum Gravity, Vol.~1, Chap.~1, 2026.

\bibitem{silvafilho2026book}
R.~M.~Silva-Filho,
\emph{Geometry, Tensors, and Quantum Gravity: The Unified Grand Synthesis Treatise},
Universidade Federal de Lavras, 2026, \href{https://doi.org/10.5281/zenodo.22290043}{DOI: 10.5281/zenodo.22290043}.

\bibitem{verstraete2010continuous}
F.~Verstraete and J.~I.~Cirac,
\emph{Continuous matrix product states for quantum fields},
Phys. Rev. Lett. \textbf{104} (2010), no.~19, 190405.

\bibitem{villani2009optimal}
C.~Villani,
\emph{Optimal Transport: Old and New},
Grundlehren der mathematischen Wissenschaften, vol.~338, Springer-Verlag, Berlin, 2009.

\bibitem{von2005transport}
M.-K.~von Renesse and K.-T.~Sturm,
\emph{Transport inequalities, gradient estimates, and Ricci curvature},
Comm. Pure Appl. Math. \textbf{58} (2005), no.~7, 923--940.

\bibitem{willmore1993riemannian}
T.~J.~Willmore,
\emph{Riemannian Geometry},
Oxford Science Publications, Oxford University Press, New York, 1993.

\end{thebibliography}

\end{document}
